\documentclass{osucourse}
\course{2137}

%% Not in the university's course data, so these come
%% from the published sheet this file was imported from.
\SetCourseField{title}{Algebra and Coordinate Geometry for Teachers}
\SetCourseField{credits}{3 credits}
\SetCourseField{semesters}{Autumn, Spring}
\SetCourseField{description}{This is one of two independent courses which follow Measurement and geome-try for teachers to provide necessary content for middle grade teachers. This course focuses on algebra, coordinate geometry, and their connections through equations in one or more unknowns. Modern and historical perspectives are woven throughout.}
\SetCourseField{prereq}{A grade of C$-$ or above in “Measurement and Geometry for Teachers” (Math 1136). A grade of C-or above in Math 1149 or 1150, or credit for 150, or math placement level L.}

\begin{document}

\CatalogDescription
\PrerequisitesAndExclusions

\begin{Purpose}
This course integrates the various types of numbers introduced in the previous course to present them as members of a single (real) number system. The notion that new numbers are discovered as solutions to equations is promoted, and motivated by connecting various equations with mathematical models.

Matrices are introduced and used as linear transformations, mainly in the plane. The complex numbers are introduced as general solutions to quadratic equations and the relationship between complex arithmetic and transformations in the plane is explored.

The course finishes with several weeks of geometry content for middle grade teachers, including more material on proofs, triangle congruence, and non-Euclidean geometry. The main example is “Taxicab geometry”, based on the l 1 norm.
\end{Purpose}

\begin{Text}
\emph{Basic Mathematics}, by Serge Lang, Springer, ISBN 9780387967875
\end{Text}

\begin{TopicsList}
\item Polynomial arithmetic as “base-x” and binomial theorem
\item Real number system
\item Polynomial equations and their roots
\item Exponential and logarithm functions
\item Complex numbers
\item Matrices
\item Complex arithmetic and linear transformations in the plane
\item Geometry proofs
\item Taxicab geometry
\end{TopicsList}

\end{document}
